\section{Limitations}
\label{sec:limitations}

Although the Gen-MOR experiment provides a useful framework for dominance-preserving generative modeling, several limitations remain. These limitations help clarify what the current experiment demonstrates and what should be improved in future work.

The dataset is generated from a synthetic MORNP-style construction process rather than from a real deployed planning system. This is useful for controlled experimentation because the dominance labels can be verified using the formal dominance relation. However, real-world planning data may contain additional structure, noise, constraints, and domain-specific dependencies that are not fully captured by the synthetic generator. Although the current 50K-plan dataset is substantially larger than the preliminary 1K dataset, it is still a controlled synthetic benchmark rather than a real-world deployment dataset.

The objectives in the current dataset are represented as generic numerical dimensions such as \texttt{obj\_1}, \texttt{obj\_2}, and so on. This makes the experiment general, but it also limits semantic interpretation. In a real planning domain, objectives may correspond to specific quantities such as cost, time, risk, energy, or safety. Future versions should connect the objective dimensions to a concrete planning environment so that generated plans can be interpreted not only structurally, but also semantically.

The current experiment focuses on AE, VAE, and CVAE models. These models are appropriate for a first Gen-MOR study because they provide a clear progression from reconstruction to probabilistic generation and class-conditioned generation. However, the experiment does not yet include more advanced generative architectures such as energy-based models, normalizing flows, diffusion models, or hybrid constraint-guided generators.

The current models are evaluated using DVR, but dominance validity is not directly optimized during training. The loss functions include reconstruction, classification, and KL-divergence terms, but they do not include a differentiable dominance penalty. As a result, a model may learn to reconstruct and classify objective vectors reasonably well without learning to generate dominance-valid samples consistently. This limitation is visible in the gap between standard reconstruction/classification metrics and DVR results.

Generated samples may vary significantly across random seeds. Even if a model achieves good average reconstruction and classification performance, its DVR may still fluctuate across runs. The multi-seed results show that DVR has much higher variance than reconstruction and classification metrics, especially for VAE and CVAE generation. This indicates that dominance-preserving generation is not yet stable. This limitation is addressed partially through multi-seed evaluation, but future models should aim to reduce this instability.

Dominance verification can become computationally expensive because it requires comparing generated samples against a verification pool under the selected objective subset. In the current 50K experiment, DVR computation is made practical using chunked vectorized verification with a reference pool of $10{,}000$ real plans and a generated-sample evaluation cap of $500$ samples. As the number of generated samples, reference plans, or objectives increases, DVR computation may become a bottleneck. Efficient dominance-checking strategies may be needed for larger-scale experiments.

The current experiment evaluates generated plans using formal dominance checks, but it does not yet test whether the generated plans improve a downstream planner, solver, or decision-support workflow. Therefore, the results demonstrate feasibility for dominance-preserving generation, but they do not yet prove real-world planning utility.